\documentclass[aspectratio=169,11pt]{beamer}

% Apresentacao de 15 minutos do Tech Challenge - Fase 3.
\usepackage[utf8]{inputenc}
\usepackage[T1]{fontenc}
\usepackage[brazil]{babel}
\usepackage{lmodern}
\usepackage{graphicx}
\usepackage{booktabs}
\usepackage{tikz}
\usetikzlibrary{positioning}
\usepackage{listings}
\usepackage{amsmath}

% Paleta inspirada na interface POS TECH: grafite, branco e magenta.
\definecolor{navy}{HTML}{090909}
\definecolor{teal}{HTML}{EC0A68}
\definecolor{mint}{HTML}{F1F1F1}
\definecolor{ink}{HTML}{E6E6E6}
\definecolor{slate}{HTML}{A5A5A5}
\definecolor{coral}{HTML}{B80852}
\definecolor{darkbg}{HTML}{151515}
\definecolor{panel}{HTML}{252525}

\usetheme{default}
\usefonttheme{professionalfonts}
\setbeamertemplate{navigation symbols}{}
\setbeamertemplate{footline}{%
  \hfill\scriptsize\color{slate}\insertshorttitle\quad|
  \quad\insertframenumber/\inserttotalframenumber\hspace{0.35cm}\vspace{0.2cm}}
\setbeamertemplate{frametitle}{%
  \vspace{0.25cm}\noindent\color{mint}\Large\bfseries\insertframetitle\par
  \vspace{0.06cm}\noindent\color{teal}\rule{\textwidth}{1.2pt}}
\setbeamercolor{background canvas}{bg=darkbg}
\setbeamercolor{normal text}{fg=ink,bg=darkbg}
\setbeamercolor{structure}{fg=teal}
\setbeamercolor{alerted text}{fg=coral}
\setbeamercolor{block title}{fg=white,bg=navy}
\setbeamercolor{block body}{fg=ink,bg=panel}
\setbeamercolor{author in head/foot}{fg=slate}

\lstset{
  basicstyle=\ttfamily\scriptsize\color{ink},
  backgroundcolor=\color{panel},
  frame=single,
  rulecolor=\color{teal!40},
  breaklines=true,
  columns=fullflexible
}

\title[Assistente Médico]{Assistente Médico com LLM, RAG e LangChain}
\subtitle{Tech Challenge -- Fase 3}
\author{Jefferson Antônio Pantoja Silva}
\date{12 de setembro de 2026}

\begin{document}

\begin{frame}[plain]
  \begin{tikzpicture}[remember picture,overlay]
    \fill[navy] (current page.south west) rectangle (current page.north east);
    \fill[teal] (current page.south east) -- ++(-4.2,0) -- ++(4.2,3.0) -- cycle;
    \fill[coral] (current page.north west) -- ++(2.4,0) -- ++(-2.4,-1.8) -- cycle;
  \end{tikzpicture}
  \vspace{1.25cm}
  {\color{white}\fontsize{28}{34}\selectfont\bfseries Assistente Médico\par}
  \vspace{0.35cm}
  {\color{mint}\Large Tech Challenge -- Fase 3\par}
  \vfill
  {\color{mint}\small Jefferson Antônio Pantoja Silva | 12 de setembro de 2026\par}
\end{frame}

\begin{frame}{O desafio}
  \begin{columns}[T,onlytextwidth]
    \column{0.54\textwidth}
    {\Large\bfseries Como responder dúvidas médicas com contexto, rastreabilidade e segurança?\par}
    \vspace{0.45cm}
    \normalsize
    O sistema deve combinar conhecimento biomédico, informações do paciente e mecanismos explícitos de proteção.
    \column{0.42\textwidth}
    \begin{block}{Objetivo}
      Construir um pipeline modular que prepare dados, ajuste uma LLM, consulte prontuários sintéticos e gere respostas contextualizadas.
    \end{block}
    \vspace{0.2cm}
    \begin{alertblock}{Limite de uso}
      Apoio informacional. Não substitui avaliação ou decisão de um profissional de saúde.
    \end{alertblock}
  \end{columns}
\end{frame}

\begin{frame}{Visão geral da solução}
  \begin{columns}[T,onlytextwidth]
    \column{0.40\textwidth}
    \vspace{0.8cm}
    \begin{block}{Ideia central}
      O modelo treinado com o dataset gera a resposta com base nas informações dos prontuários sintéticos criados a partir do dataset
    \end{block}
    \vspace{0.2cm}
    \small\color{slate} O pipeline separa preparação, conhecimento do modelo e contexto do paciente.
    \column{0.57\textwidth}
    \centering
    \resizebox{\linewidth}{!}{%
    \begin{tikzpicture}[every node/.style={font=\tiny},
      box/.style={draw=teal, fill=panel, text=ink, rounded corners, minimum height=0.52cm, text width=3.1cm, align=center},
      topbox/.style={draw=teal, fill=panel, text=ink, rounded corners, minimum height=0.52cm, text width=3.6cm, align=center},
      arrow/.style={->, very thick, color=teal}]
      \node[topbox] (sources) at (0,2.10) {MedQuAD + PubMedQA};
      \node[topbox] (curation) at (0,1.45) {Leitura + curadoria};
      \node[box] (records) at (-1.75,0.78) {Prontuários sintéticos};
      \node[box] (dataset) at (1.75,0.78) {Dataset JSONL};
      \node[box] (normalize) at (-1.75,0.10) {Normalização + JSONL};
      \node[box] (train) at (1.75,0.10) {QLoRA + SFTTrainer};
      \node[box] (embed) at (-1.75,-0.58) {Embeddings + FAISS};
      \node[box] (model) at (1.75,-0.58) {fine-tuned model};
      \node[box] (lang) at (0,-1.26) {LangChain + LangGraph};
      \node[box] (answer) at (0,-1.94) {Resposta contextual};
      \draw[arrow] (sources) -- (curation);
      \draw[arrow] (curation) -- (records);
      \draw[arrow] (curation) -- (dataset);
      \draw[arrow] (records) -- (normalize);
      \draw[arrow] (normalize) -- (embed);
      \draw[arrow] (dataset) -- (train);
      \draw[arrow] (train) -- (model);
      \draw[arrow] (model) -- (lang);
      \draw[arrow] (embed) -- (lang);
      \draw[arrow] (lang) -- (answer);
    \end{tikzpicture}
    }
  \end{columns}
\end{frame}

\begin{frame}{Dados: curadoria antes do modelo}
  \small
  \begin{columns}[T,onlytextwidth]
    \column{0.42\textwidth}
    \textbf{Fontes públicas}
    \begin{itemize}
      \item MedQuAD: XML com perguntas e respostas do NIH.
      \item PubMedQA: JSON com pergunta, contexto e resposta biomédica.
    \end{itemize}
    \column{0.54\textwidth}
    \begin{block}{Curadoria}
      Normalização, validação, remoção de vazios e duplicidades. O campo \texttt{source} é preservado para rastreabilidade e o JSONL é escrito incrementalmente.
    \end{block}
    \vspace{0.05cm}
  \end{columns}
  \vspace{0.08cm}
  \begin{block}{Saída: \texttt{finetuning\_qa.jsonl}}
      \scriptsize
      \texttt{\{ "source": "MedQuAD:arquivo:pid",}\\[-0.02cm]
      \texttt{  "text": "ANSWER THE QUESTION.}\\
      \texttt{[|Context|] ...[|eContext|]}\\[-0.02cm]
      \texttt{[|Question|] ...[|eQuestion|]}\\[-0.02cm]
      \texttt{[|Answer|] ...[|eAnswer|]" \}}\\[0.08cm]
      \textbf{Mesmo template na inferência:} reduz o desalinhamento entre treino e uso do modelo.
  \end{block}
\end{frame}

\begin{frame}{Fine-tuning no Colab}
  \begin{columns}[T,onlytextwidth]
    \column{0.55\textwidth}
    \begin{enumerate}
      \item Dataset separado em treino e teste.
      \item Modelo base \texttt{Llama-3.2-1B-Instruct} em 4-bit.
      \item Adaptação com QLoRA e Unsloth.
      \item Treino supervisionado com \texttt{trl.SFTTrainer}.
      \item Merge do adapter e salvamento do modelo final.
    \end{enumerate}
    \column{0.4\textwidth}
    \begin{block}{Parâmetros principais}
      \small
      LoRA $r=16$ | alpha $=32$\\
      dropout $=0{,}05$\\
      batch $=2$ | acumulação $=8$\\
      learning rate $=2\times10^{-4}$\\
      1 época | BF16/FP16
    \end{block}
  \end{columns}
  \vspace{0.25cm}
  \centering\small O treinamento é executado no Google Colab; a inferência usa o modelo local mesclado.
\end{frame}

\begin{frame}{Prontuários sintéticos e rastreabilidade}
  \begin{columns}[T,onlytextwidth]
    \column{0.5\textwidth}
    Um caso curado é convertido em um prontuário fictício com estrutura consistente:
    \begin{itemize}
      \item queixa principal e histórico;
      \item medicamentos e sinais vitais;
      \item avaliação e plano;
      \item origem e identificador do paciente.
      \item OpenAI API (gpt-4o-mini)
    \end{itemize}
    \column{0.46\textwidth}
    \begin{block}{Confiabilidade do pipeline}
      Os lotes são validados pela presença de exatamente uma ocorrência de cada \texttt{source}. Respostas inválidas podem ser reprocessadas até três vezes.
    \end{block}
    \begin{block}{Retomada}
      O JSONL é gravado progressivamente e checkpoints evitam repetir registros já concluídos.
    \end{block}
  \end{columns}
\end{frame}

\begin{frame}{Avaliação do modelo fine-tuned}
  \small
  \begin{columns}[T,onlytextwidth]
    \column{0.46\textwidth}
    \textbf{Protocolo}
    \begin{itemize}
      \item Amostra aleatória de 500 itens do teste.
      \item Teste separado via \texttt{train\_test\_split}.
      \item Apenas o fine-tuned; sem comparação com o modelo base.
    \end{itemize}
    \vspace{0.15cm}
    \begin{alertblock}{Interpretação}
      Métricas textuais não comprovam factualidade ou validade clínica.
    \end{alertblock}
    \column{0.48\textwidth}
    \begin{block}{Resultados registrados}
      \centering
      \begin{tabular}{lr}
        \toprule
        \textbf{Métrica} & \textbf{Resultado}\\
        \midrule
        Exact Match & 0,0000\\
        Token F1 & 0,3280\\
        ROUGE-L & 0,2717\\
        BERTScore F1 & 0,8354\\
        \bottomrule
      \end{tabular}
    \end{block}
    \vspace{0.15cm}
    \footnotesize Exact Match nulo e Token F1/ROUGE-L moderados indicam baixa correspondência literal. O BERTScore sugere similaridade semântica, mas não garante correção factual.
  \end{columns}
\end{frame}

\begin{frame}{Contexto clínico com controle de risco}
  \begin{columns}[T,onlytextwidth]
    \column{0.48\textwidth}
    \begin{block}{Recuperação}
      \begin{itemize}
        \item Busca semântica com \texttt{all-MiniLM-L6-v2}.
        \item Até quatro documentos relevantes.
        \item Filtro opcional pelo \texttt{patient\_id}.
        \item Fontes retornam junto da resposta.
      \end{itemize}
    \end{block}
    \column{0.48\textwidth}
    \begin{alertblock}{Salvaguardas}
      \begin{itemize}
        \item O campo \texttt{plan} fica fora do contexto RAG.
        \item Pedidos de medicamento ou tratamento são revisados.
        \item A resposta inclui aviso de avaliação profissional.
        \item Há fallback se a revisão externa falhar.
      \end{itemize}
    \end{alertblock}
  \end{columns}
\end{frame}

\begin{frame}{Arquitetura e fluxo RAG}
  \begin{columns}[T,onlytextwidth]
    \column{0.37\textwidth}
    \textbf{Clean Architecture}
    \begin{itemize}
      \item \texttt{domain}: modelos imutáveis.
      \item \texttt{application}: casos de uso e portas.
      \item \texttt{infrastructure}: LLM, FAISS, OpenAI e Langfuse.
      \item \texttt{presentation}: controladores CLI.
    \end{itemize}
    \column{0.59\textwidth}
    \centering
    \begin{tikzpicture}[node distance=0.12cm, every node/.style={font=\scriptsize},
      flow/.style={draw=teal, fill=panel, text=ink, rounded corners, minimum width=5.2cm, minimum height=0.48cm, align=center},
      arr/.style={->, color=navy}]
      \node[flow] (t) {translate\_question};
      \node[flow, below=of t] (r) {retrieve: embeddings + FAISS};
      \node[flow, below=of r] (c) {build\_context: registros relevantes};
      \node[flow, below=of c] (g) {generate: LLM local fine-tuned};
      \node[flow, below=of g] (a) {translate\_answer};
      \node[flow, below=of a] (v) {review\_output + fontes};
      \foreach \x/\y in {t/r,r/c,c/g,g/a,a/v}{\draw[arr] (\x) -- (\y);}
    \end{tikzpicture}
  \end{columns}
\end{frame}

\begin{frame}[plain]
  \centering
  \vfill
  {\color{teal}\fontsize{32}{38}\selectfont\bfseries Demonstração!\par}
  \vfill
\end{frame}

\end{document}
